
\documentclass{article}
\usepackage{colortbl}
\usepackage{makecell}
\usepackage{multirow}
\usepackage{supertabular}

\begin{document}

\newcounter{utterance}

\centering \large Interaction Transcript for game `cladder', experiment `full\_v1.5\_default', episode 3375 with score{-}Qwen3.5{-}9B{-}sp.
\vspace{24pt}

{ \footnotesize  \setcounter{utterance}{1}
\setlength{\tabcolsep}{0pt}
\begin{supertabular}{c@{$\;$}|p{.15\linewidth}@{}p{.15\linewidth}p{.15\linewidth}p{.15\linewidth}p{.15\linewidth}p{.15\linewidth}}
    \# & \multicolumn{2}{c}{Player} && \multicolumn{2}{c}{Game Master} \\
    \hline

    \theutterance \stepcounter{utterance}  
    & & & \multicolumn{4}{p{0.6\linewidth}}{
        \cellcolor[rgb]{0.9,0.9,0.9}{
            \makecell[{{p{\linewidth}}}]{
                \texttt{\tiny{[P1$\langle$GM]}}
                \texttt{You are an expert in causal inference. The following question is not a typical commonsense query, but rather a meticulously designed question created by a professor specializing in causal inference, intended to assess the students' mastery of the course content.} \\
\texttt{Imagine a self{-}contained, hypothetical world with only the following conditions, and without any unmentioned factors or causal relationships: Full moon has a direct effect on sprinkler and weather. Weather has a direct effect on ground. Sprinkler has a direct effect on ground. The overall probability of full moon is 31\%. The probability of no full moon and wet ground is 33\%. The probability of full moon and wet ground is 19\%. Is the chance of wet ground smaller when observing full moon?} \\
\texttt{Start your answer with 'Yes' or 'No', followed by additional reasoning or evidence to support your explanation.} \\
            }
        }
    }
     \\ \\

    \theutterance \stepcounter{utterance}  
    & \multicolumn{4}{p{0.6\linewidth}}{
        \cellcolor[rgb]{0.561,0.737,0.561}{
            \makecell[{{p{\linewidth}}}]{
                \texttt{\tiny{[P1$\rangle$GM]}}
                \texttt{No} \\
\\ 
\texttt{**Reasoning:**} \\
\\ 
\texttt{To determine if the chance of wet ground is smaller when observing a full moon, we must compare the conditional probability \$P(\textbackslash{}text\{Wet\} \textbackslash{}mid \textbackslash{}text\{Full Moon\})\$ with the marginal probability \$P(\textbackslash{}text\{Wet\})\$ (or at least compare it to the alternative scenario, \$P(\textbackslash{}text\{Wet\} \textbackslash{}mid \textbackslash{}text\{No Full Moon\})\$).} \\
\\ 
\texttt{Based on the provided data:} \\
\texttt{1.  **Calculate the marginal probability of wet ground (\$P(\textbackslash{}text\{Wet\})\$):**} \\
\texttt{    The problem states the probability of "no full moon and wet ground" is 33\% and "full moon and wet ground" is 19\%. Assuming these are the only two mutually exclusive states for the ground condition relative to the moon phase:} \\
\texttt{    \$\$P(\textbackslash{}text\{Wet\}) = P(\textbackslash{}text\{Wet\} \textbackslash{}cap \textbackslash{}text\{No Full Moon\}) + P(\textbackslash{}text\{Wet\} \textbackslash{}cap \textbackslash{}text\{Full Moon\})\$\$} \\
\texttt{    \$\$P(\textbackslash{}text\{Wet\}) = 33\textbackslash{}\% + 19\textbackslash{}\% = 52\textbackslash{}\%\$\$} \\
\\ 
\texttt{2.  **Calculate the conditional probability of wet ground given a full moon (\$P(\textbackslash{}text\{Wet\} \textbackslash{}mid \textbackslash{}text\{Full Moon\})\$):**} \\
\texttt{    Using the definition of conditional probability:} \\
\texttt{    \$\$P(\textbackslash{}text\{Wet\} \textbackslash{}mid \textbackslash{}} \\
            }
        }
    }
    & & \\ \\

    \theutterance \stepcounter{utterance}  
    & & & \multicolumn{2}{p{0.3\linewidth}}{
        \cellcolor[rgb]{0.9,0.9,0.9}{
            \makecell[{{p{\linewidth}}}]{
                \texttt{\tiny{[GM$|$GM]}}
                \texttt{No} \\
            }
        }
    }
    & & \\ \\

    \theutterance \stepcounter{utterance}  
    & & & \multicolumn{2}{p{0.3\linewidth}}{
        \cellcolor[rgb]{0.9,0.9,0.9}{
            \makecell[{{p{\linewidth}}}]{
                \texttt{\tiny{[GM$|$GM]}}
                \texttt{no} \\
            }
        }
    }
    & & \\ \\

    \theutterance \stepcounter{utterance}  
    & & & \multicolumn{2}{p{0.3\linewidth}}{
        \cellcolor[rgb]{0.9,0.9,0.9}{
            \makecell[{{p{\linewidth}}}]{
                \texttt{\tiny{[GM$|$GM]}}
                \texttt{game\_result = WIN} \\
            }
        }
    }
    & & \\ \\

\end{supertabular}
}

\end{document}
